\documentclass[11pt,a4paper]{article}
\usepackage[utf8]{inputenc}
\usepackage[T1]{fontenc}
\usepackage{amsmath,amsfonts,amssymb}
\usepackage{graphicx}
\usepackage{hyperref}
\usepackage{listings}
\usepackage{color}
\usepackage{booktabs}
\usepackage{float}
\usepackage{algorithm}
\usepackage{algorithmic}
\usepackage{geometry}
\geometry{margin=1in}
\definecolor{zenblue}{RGB}{41,121,255}
\hypersetup{colorlinks=true,linkcolor=zenblue,urlcolor=zenblue,citecolor=zenblue}

\title{\textbf{Zen Alignment: RLHF and Constitutional AI at Scale}\\
\large Technical Report v2025.04}
\author{Zen LM Research Team\\
\texttt{research@zenlm.org}}
\date{April 2025}

\begin{document}
\maketitle

\begin{abstract}
We describe the alignment infrastructure for the Zen model family, covering reward
modeling, proximal policy optimization, and a constitutional AI framework that reduces
dependence on human preference annotation. Key contributions include multi-reward
fusion that jointly optimizes helpfulness, harmlessness, and honesty; an adaptive
KL penalty that prevents mode collapse while maintaining policy expressiveness; and
a synthetic preference generation pipeline that scales annotation by $12\times$ at
$\sim$90\% of human annotator agreement. On head-to-head human evaluation, aligned
Zen models achieve 68.4\% win rate against SFT-only baselines, with harmlessness
improving from 71.2\% to 94.7\% and helpfulness from 76.8\% to 89.3\%.
\end{abstract}

\section{Introduction}

Aligning large language models to human values is a multi-objective optimization problem:
models must be simultaneously \textit{helpful} (completing user tasks accurately),
\textit{harmless} (avoiding dangerous or offensive outputs), and \textit{honest}
(refusing to assert false claims). These objectives conflict in non-trivial ways.
A model optimized purely for helpfulness tends to hallucinate; one optimized purely
for harmlessness becomes overly cautious and refuses legitimate requests.

The Zen alignment framework addresses this tension through three innovations:

\begin{enumerate}
  \item \textbf{Multi-reward fusion} with learned weighting across reward dimensions
  \item \textbf{Adaptive KL penalty} that adjusts dynamically during training
  \item \textbf{Constitutional synthetic preference generation} at scale
\end{enumerate}

We detail each component, report experimental results, and analyze the
helpfulness-harmlessness tradeoff surface achieved by our approach.

\section{Background}

\subsection{Reinforcement Learning from Human Feedback}

RLHF~\cite{christiano2017deep,ouyang2022rlhf} trains a reward model $r_\phi$ on
human preference data, then uses reinforcement learning (typically PPO~\cite{schulman2017ppo})
to maximize expected reward subject to a KL constraint against the reference policy.

The standard objective is:
\begin{equation}
J(\theta) = \mathbb{E}_{x \sim \mathcal{D}, y \sim \pi_\theta(\cdot|x)}\left[r_\phi(x,y) - \beta \cdot \text{KL}[\pi_\theta(\cdot|x) \| \pi_\text{ref}(\cdot|x)]\right]
\end{equation}

A fixed $\beta$ creates instability: too small permits reward hacking; too large
prevents meaningful policy improvement. We address this with adaptive scheduling.

\subsection{Constitutional AI}

Constitutional AI~\cite{bai2022constitutional} replaces human preference labels
with model-generated critiques based on a set of principles (the ``constitution'').
This scales annotation but risks distribution shift when the critiquing model shares
failure modes with the model being trained. Zen addresses this via an independent
critique model trained on a held-out safety dataset.

\section{Reward Modeling}

\subsection{Multi-Reward Architecture}

Rather than a single scalar reward, we train three specialized reward heads:
\begin{align}
r_\text{help}(x,y) &: \text{Task completion quality and accuracy} \\
r_\text{harm}(x,y) &: \text{Safety and policy compliance} \\
r_\text{honest}(x,y) &: \text{Factual accuracy and calibration}
\end{align}

All three heads share a Zen-7B backbone with task-specific linear projection heads
of dimension 1 (scalar reward). The backbone is initialized from the SFT checkpoint
to preserve instruction-following representations.

\subsubsection{Training Data}

\begin{table}[H]
\centering
\begin{tabular}{lrrr}
\toprule
\textbf{Reward Dimension} & \textbf{Human Pairs} & \textbf{Synthetic Pairs} & \textbf{Total} \\
\midrule
Helpfulness & 120K & 980K & 1.1M \\
Harmlessness & 80K & 420K & 500K \\
Honesty & 60K & 340K & 400K \\
\midrule
Total & 260K & 1.74M & 2.0M \\
\bottomrule
\end{tabular}
\caption{Reward model training data composition.}
\end{table}

Human annotators rated response pairs on a 4-point scale per dimension.
Synthetic pairs were generated via constitutional critiques (Section 4) and filtered
to pairs where the constitutional model assigned high confidence ($>0.85$).

\subsubsection{Loss Function}

Multi-reward training minimizes:
\begin{equation}
\mathcal{L}_\text{RM} = \sum_{d \in \{\text{help,harm,honest}\}} \lambda_d \cdot \mathbb{E}\left[-\log \sigma\left(r_\phi^d(x, y_w) - r_\phi^d(x, y_l)\right)\right]
\end{equation}

with $\lambda_\text{help} = 0.45$, $\lambda_\text{harm} = 0.35$, $\lambda_\text{honest} = 0.20$.
These weights are set via a small held-out human evaluation of $5000$ pairs.

\subsection{Reward Model Accuracy}

\begin{table}[H]
\centering
\begin{tabular}{lcccc}
\toprule
\textbf{RM Variant} & \textbf{Help Acc.} & \textbf{Harm Acc.} & \textbf{Honest Acc.} & \textbf{Avg.} \\
\midrule
Single scalar & 72.1\% & 78.4\% & 68.3\% & 72.9\% \\
Multi-head (ours) & \textbf{79.8\%} & \textbf{85.2\%} & \textbf{76.4\%} & \textbf{80.5\%} \\
Multi-head + synth & 78.9\% & 84.7\% & 76.1\% & 79.9\% \\
\bottomrule
\end{tabular}
\caption{Reward model accuracy on held-out human preference pairs (6K pairs per dimension).}
\end{table}

The multi-head architecture improves over a single scalar reward model by 7.6 points
average accuracy. Adding synthetic data slightly reduces accuracy on human pairs
(distribution shift) but enables higher-quality policy training.

\section{Constitutional Synthetic Preference Generation}

\subsection{Constitution Definition}

The Zen constitution contains 42 principles organized into five categories:

\begin{enumerate}
  \item \textbf{Factual accuracy} (10 principles): truthfulness, uncertainty expression,
    source attribution, numerical accuracy
  \item \textbf{Safety} (12 principles): refusal of harmful requests, child protection,
    no dangerous instructions, bias avoidance
  \item \textbf{Privacy} (6 principles): no PII disclosure, data minimization
  \item \textbf{Helpfulness} (8 principles): task relevance, completeness, clarity
  \item \textbf{Fairness} (6 principles): demographic balance, representation
\end{enumerate}

\subsection{Synthetic Preference Generation Algorithm}

\begin{algorithm}[H]
\caption{Constitutional Synthetic Preference Generation}
\begin{algorithmic}[1]
\REQUIRE Prompt set $\mathcal{P}$, constitution $\mathcal{C}$, model $\pi$, critique model $\pi_c$
\ENSURE Synthetic preference pairs $\mathcal{S}$
\STATE $\mathcal{S} \leftarrow \emptyset$
\FOR{each prompt $p \in \mathcal{P}$}
  \STATE $y_0 \leftarrow \pi(p)$ \COMMENT{Initial response}
  \STATE Sample $c_i \sim \mathcal{C}$ uniformly
  \STATE $\text{critique} \leftarrow \pi_c(p, y_0, c_i)$ \COMMENT{Generate critique}
  \STATE $y_1 \leftarrow \pi(p, \text{critique})$ \COMMENT{Revised response}
  \STATE $\text{conf} \leftarrow \pi_c.\text{score}(p, y_0, y_1, c_i)$ \COMMENT{Confidence}
  \IF{$\text{conf} > 0.85$}
    \STATE $\mathcal{S} \leftarrow \mathcal{S} \cup \{(p, y_1, y_0)\}$ \COMMENT{Revised preferred}
  \ENDIF
\ENDFOR
\RETURN $\mathcal{S}$
\end{algorithmic}
\end{algorithm}

The critique model $\pi_c$ is a Zen-32B model fine-tuned on 80K human-written
critique-revision pairs, held separate from the policy being aligned to prevent
feedback loops.

\subsection{Synthetic vs. Human Preference Agreement}

\begin{table}[H]
\centering
\begin{tabular}{lcc}
\toprule
\textbf{Principle Category} & \textbf{Synth-Human Agreement} & \textbf{Human-Human Agreement} \\
\midrule
Factual accuracy & 88.2\% & 91.4\% \\
Safety & 91.3\% & 93.7\% \\
Privacy & 87.6\% & 89.1\% \\
Helpfulness & 83.4\% & 87.2\% \\
Fairness & 79.8\% & 84.3\% \\
\midrule
Average & \textbf{86.1\%} & \textbf{89.1\%} \\
\bottomrule
\end{tabular}
\caption{Agreement between synthetic preferences and human annotators vs. inter-annotator agreement.}
\end{table}

\section{PPO Training with Adaptive KL}

\subsection{Fused Reward Signal}

The fused reward used during PPO combines the three reward dimensions:
\begin{equation}
r_\text{fused}(x,y) = \mathbf{w}^\top \mathbf{r}(x,y) - \beta_t \cdot \text{KL}[\pi_\theta(\cdot|x)\|\pi_\text{ref}(\cdot|x)]
\end{equation}

where $\mathbf{w} = (w_\text{help}, w_\text{harm}, w_\text{honest})$ is a learned
weight vector and $\beta_t$ is the adaptive KL penalty.

\subsection{Adaptive KL Penalty}

We introduce a PID-style adaptive KL controller:

\begin{equation}
\beta_{t+1} = \beta_t \cdot \exp\left(\kappa \cdot \left(\overline{\text{KL}}_t - \text{KL}_\text{target}\right)\right)
\end{equation}

where $\overline{\text{KL}}_t$ is the exponential moving average of KL divergence,
$\text{KL}_\text{target} = 0.1$ nats is the target, and $\kappa = 0.1$ is the
controller gain. $\beta$ is clipped to $[0.01, 1.0]$.

\begin{algorithm}[H]
\caption{PPO with Adaptive KL and Multi-Reward Fusion}
\begin{algorithmic}[1]
\REQUIRE Policy $\pi_\theta$, reference $\pi_\text{ref}$, reward heads $r^d_\phi$, weights $\mathbf{w}$
\STATE Initialize $\beta_0 = 0.04$, $\overline{\text{KL}}_0 = 0$
\FOR{step $t = 1 \ldots T$}
  \STATE Sample prompts $\{x_i\}$ from dataset
  \STATE Generate responses $\{y_i\} \sim \pi_\theta(\cdot|x_i)$
  \STATE Compute $r^d_i = r^d_\phi(x_i, y_i)$ for each dimension $d$
  \STATE $r_i = \mathbf{w}^\top \mathbf{r}_i - \beta_t \cdot \text{KL}_i$
  \STATE Compute advantages $\hat{A}_i$ via GAE with value baseline
  \FOR{$k = 1 \ldots K_\text{epochs}$}
    \STATE Update $\theta$ via clipped PPO objective: $\min(r_t \hat{A}, \text{clip}(r_t, 1\pm\epsilon)\hat{A})$
  \ENDFOR
  \STATE $\overline{\text{KL}}_{t+1} \leftarrow 0.9 \cdot \overline{\text{KL}}_t + 0.1 \cdot \overline{\text{KL}}_{\text{batch}}$
  \STATE $\beta_{t+1} \leftarrow \text{clip}\left(\beta_t \exp(\kappa(\overline{\text{KL}}_t - 0.1)), 0.01, 1.0\right)$
\ENDFOR
\end{algorithmic}
\end{algorithm}

\subsection{Policy Optimization Metrics}

\begin{table}[H]
\centering
\begin{tabular}{lcccc}
\toprule
\textbf{Method} & \textbf{Reward} & \textbf{KL (nats)} & \textbf{Win Rate vs SFT} & \textbf{Refusal Rate} \\
\midrule
SFT baseline & -- & -- & 50.0\% & 8.2\% \\
PPO fixed $\beta=0.01$ & 0.83 & 0.41 & 61.2\% & 18.3\% \\
PPO fixed $\beta=0.10$ & 0.71 & 0.08 & 57.4\% & 11.4\% \\
PPO adaptive (ours) & \textbf{0.89} & \textbf{0.10} & \textbf{68.4\%} & \textbf{13.7\%} \\
\bottomrule
\end{tabular}
\caption{PPO training results on Zen-7B. Win rate from 2K human-judged comparisons.}
\end{table}

\section{Helpfulness-Harmlessness Tradeoff}

A key concern in alignment is the tradeoff between helpfulness and harmlessness:
models that refuse more often are safer but less useful. We characterize this frontier
by sweeping $w_\text{harm}$ from 0.1 to 0.9 while normalizing other weights:

\begin{table}[H]
\centering
\begin{tabular}{cccc}
\toprule
\textbf{$w_\text{harm}$} & \textbf{Helpfulness (\%)} & \textbf{Harmlessness (\%)} & \textbf{Over-refusal (\%)} \\
\midrule
0.1 & 91.2 & 78.3 & 4.1 \\
0.2 & 90.4 & 84.7 & 6.3 \\
0.35 (default) & 89.3 & 94.7 & 13.7 \\
0.5 & 85.1 & 97.1 & 22.4 \\
0.7 & 76.8 & 99.1 & 38.2 \\
\bottomrule
\end{tabular}
\caption{Pareto frontier of helpfulness vs. harmlessness at varying harm reward weight.}
\end{table}

Our default setting ($w_\text{harm}=0.35$) achieves an over-refusal rate of 13.7\%
(vs. 8.2\% for unaligned SFT) while improving harmlessness from 71.2\% to 94.7\%.

\section{Human Preference Evaluation}

\subsection{Evaluation Protocol}

We conducted a human preference study with 48 annotators evaluating 3000 prompt-response
pairs per model comparison. Annotators rated each response on helpfulness,
harmlessness, and honesty on a 5-point scale, and made an overall preference judgment.

\subsection{Results}

\begin{table}[H]
\centering
\begin{tabular}{lcccc}
\toprule
\textbf{Comparison} & \textbf{Win} & \textbf{Tie} & \textbf{Loss} & \textbf{Net Preference} \\
\midrule
Zen-7B-Aligned vs. SFT & 68.4\% & 12.3\% & 19.3\% & +49.1\% \\
Zen-32B-Aligned vs. SFT & 71.2\% & 11.8\% & 17.0\% & +54.2\% \\
Zen-7B single-RM vs. SFT & 59.8\% & 14.1\% & 26.1\% & +33.7\% \\
\bottomrule
\end{tabular}
\caption{Human preference win rates. Evaluations conducted by independent annotators.}
\end{table}

\section{Analysis}

\subsection{Reward Hacking Behavior}

We observe reward hacking manifests differently per reward dimension:
\begin{itemize}
  \item \textbf{Helpfulness}: Length exploitation (longer responses score higher, model
    learns to pad); mitigated by length-normalized reward.
  \item \textbf{Harmlessness}: Excessive hedging (adding disclaimers to benign responses);
    mitigated by over-refusal penalty term.
  \item \textbf{Honesty}: Uncertainty inflation (saying ``I'm not sure'' to avoid
    factual claims); mitigated by calibration reward signal.
\end{itemize}

\subsection{Constitutional Training Contribution}

Ablating constitutional training (Section 3 of zen-training-methodology) shows
12-point harmlessness degradation at the same refusal rate, confirming that
constitutional training instills genuine safety understanding rather than
pattern-matching refusal.

\subsection{Scaling Reward Model Size}

Larger reward models yield monotonically better policy quality:

\begin{table}[H]
\centering
\begin{tabular}{lcc}
\toprule
\textbf{RM Size} & \textbf{Win Rate vs. SFT} & \textbf{Human-RM Correlation} \\
\midrule
1.3B & 61.2\% & 0.74 \\
7B (default) & 68.4\% & 0.81 \\
32B & 71.8\% & 0.85 \\
\bottomrule
\end{tabular}
\caption{Effect of reward model size on policy quality.}
\end{table}

Diminishing returns above 7B RM size suggest a ceiling in preference prediction
accuracy given our annotation methodology.

\section{Conclusion}

The Zen alignment framework demonstrates that multi-objective reward modeling with
adaptive KL control and constitutional synthetic preference generation produces
models that substantially improve over SFT baselines on all alignment dimensions.
The 68.4\% human preference win rate and 23.5-point harmlessness improvement
(71.2\% $\to$ 94.7\%) validate the framework at the 7B scale, with consistent
gains at 32B.

Future directions include online RLHF where the reward model is updated in parallel
with the policy, debate-based preference elicitation for complex reasoning tasks,
and multi-lingual alignment extending the constitutional framework to non-English
principles.

\bibliographystyle{plain}
\begin{thebibliography}{99}
\bibitem{christiano2017deep} Christiano et al. (2017). Deep Reinforcement Learning from Human Preferences. \textit{NeurIPS}.
\bibitem{ouyang2022rlhf} Ouyang et al. (2022). Training language models to follow instructions with human feedback. \textit{NeurIPS}.
\bibitem{schulman2017ppo} Schulman et al. (2017). Proximal Policy Optimization Algorithms. \textit{arXiv:1707.06347}.
\bibitem{bai2022constitutional} Bai et al. (2022). Constitutional AI: Harmlessness from AI Feedback. \textit{arXiv:2212.08073}.
\bibitem{stiennon2020rlhf} Stiennon et al. (2020). Learning to summarize from human feedback. \textit{NeurIPS}.
\end{thebibliography}

\end{document}
